\documentclass[11pt, letterpaper]{article}
\input{../../homework.tex}
\usepackage{titlepageBU}
\title{Problem Set 3}
\courseID{PY 451}
\courseSection{A1}
\begin{document}
\makeproblem
\section{Representations and Basis States}
Even though our space is infinite-dimensional, we can isolate $v_1,v_2$ to consider a 2-dimensional subspace. The most general possible choice in this subspace, as illustrated by equation (1.1), is
\[
	\ket{1}=\cos \theta e^{i\phi_1} v_1+\sin \theta e^{i\phi_2} v_2,
\]
\[
	\ket{2}=\sin \theta e^{i\phi_3} v_1-\cos \theta e^{i\left( \phi_3+\phi_1-\phi_2 \right) } v_2
.\]

Idfk what to do for b
\section{Operators}
\begin{enumerate}
	\item h
	\item Note that $h(x)g(x)=0$ (except possibly on $\left\{ 1/3,2/3 \right\} $ depending on how the drawing is interpreted, but this set has lebesgue measure 0 so the following conclusion still holds). It follows that
		\[
			\int_{0}^{1} h(y)g(y)g(x) \,\mathrm{d} y=0
		.\]
		Therefore,
		\[
			(Ah)(x)=\int_{0}^{1} h(y)^2h(x) \,\mathrm{d} y+\int_{0}^{1} h(y)v(y)v(x) \,\mathrm{d} x
		.\]
		Note that $v(x)=\frac{1}{3} h(x)$. It follows that
		\[
			(Ah)(x)=\int_{0}^{1} h(y)^2h(x) \,\mathrm{d} y+\int_{0}^{1} \frac{1}{9} h(y)^2 h(x)\,\mathrm{d} y=\frac{10}{9} h(x)\int_{0}^{1} h(y)^2 \,\mathrm{d} y
		.\]
		Also note that since $0^2=0$ and $1^2=1$, we have that $h(x)^2=h(x)$, so
		\[
			(Ah)(x)=\frac{10}{9} h(x)\int_{0}^{1} h(y) \,\mathrm{d} y=\frac{20}{27} h(x) 
		.\]
	\item bruh
\end{enumerate}

\end{document}
